\documentclass[11pt]{article}

\usepackage[T1]{fontenc}
\usepackage[utf8]{inputenc}
\usepackage{lmodern}
\usepackage[margin=1.06in]{geometry}
\usepackage{amsmath,amssymb,mathtools,amsthm}
\usepackage{aliascnt}
\usepackage{mathrsfs}
\usepackage{microtype}
\usepackage{enumitem}
\usepackage{xcolor}
\usepackage{url}
\usepackage{hyperref}

\hypersetup{
  colorlinks=true,
  linkcolor=blue!45!black,
  citecolor=green!35!black,
  urlcolor=blue!55!black,
  pdftitle={Algebraic Exponential Integrals: Transcendence and Linear Independence},
  pdfauthor={Christopher D. Long},
  pdfsubject={E-functions, exponential integrals, transcendence, linear independence},
  pdfkeywords={E-function, Siegel--Shidlovskii theorem, exponential integral, transcendence, linear independence, polynomial moments}
}

\numberwithin{equation}{section}
\setlength{\emergencystretch}{3em}

\newtheorem{theorem}{Theorem}[section]

\newaliascnt{proposition}{theorem}
\newtheorem{proposition}[proposition]{Proposition}
\aliascntresetthe{proposition}

\newaliascnt{lemma}{theorem}
\newtheorem{lemma}[lemma]{Lemma}
\aliascntresetthe{lemma}

\newaliascnt{corollary}{theorem}
\newtheorem{corollary}[corollary]{Corollary}
\aliascntresetthe{corollary}

\theoremstyle{definition}
\newaliascnt{definition}{theorem}
\newtheorem{definition}[definition]{Definition}
\aliascntresetthe{definition}

\theoremstyle{remark}
\newaliascnt{remark}{theorem}
\newtheorem{remark}[remark]{Remark}
\aliascntresetthe{remark}

\usepackage[capitalize,noabbrev]{cleveref}

\newcommand{\C}{\mathbb C}
\newcommand{\Q}{\mathbb Q}
\newcommand{\Qbar}{\overline{\mathbb Q}}
\newcommand{\Aone}{\mathbb A^1}
\newcommand{\Span}{\operatorname{span}}
\newcommand{\Tr}{\operatorname{Tr}}
\newcommand{\Spec}{\operatorname{Spec}}
\newcommand{\End}{\operatorname{End}}
\newcommand{\ord}{\operatorname{ord}}
\newcommand{\diag}{\operatorname{diag}}
\newcommand{\Mcal}{\mathcal M}
\newcommand{\Ecal}{\mathcal E}
\newcommand{\Rcal}{\mathcal R}
\newcommand{\Ncal}{\mathcal N}
\newcommand{\Ocal}{\mathcal O}
\newcommand{\Dmod}{\mathcal D}
\newcommand{\Lcal}{\mathscr L}
\newcommand{\Vcal}{\mathscr V}
\newcommand{\dd}{\,d}

\title{Algebraic Exponential Integrals:\\
Transcendence and Linear Independence}
\author{Christopher D. Long}
\date{August 2026}

\begin{document}

\maketitle

\begin{center}
\small
Headlamp Software, Lexington, Kentucky, USA\\
\texttt{galizur@gmail.com}
\end{center}

\begin{abstract}
Let $q\in\Qbar[x]$ be nonconstant, let $\xi\in\Qbar^\times$, and let
$\Delta=\sum_\alpha d_\alpha[\alpha]$ be a finite degree-zero zero-chain with
$d_\alpha,\alpha\in\Qbar$.  The entire function
\[
  I_{q,\Delta}(z)=\int_\Delta e^{zq(x)}\,dx
\]
is defined by integration over any one-chain with boundary $\Delta$.  We prove
the algebraic-value dichotomy
\[
  I_{q,\Delta}(\xi)\in\Qbar
  \quad\Longleftrightarrow\quad
  I_{q,\Delta}(z)\equiv0.
\]
Thus every nonzero algebraic-parameter value of an algebraic exponential
integral is transcendental.  More generally, all $\Qbar$-linear relations among
finitely many such values, together with $1$, are exactly the specializations
of their functional relations.  This yields linear independence, for example,
for exponential integrals from a common endpoint to algebraic points whose
$q$-values are pairwise distinct regular values.

For $\deg q\ge2$, the proof uses an explicit inhomogeneous differential system
for the moment functions
$\int_\Delta x^j e^{zq(x)}\,dx$.  A complete-reducibility theorem for its
homogeneous part, combined with an elementary local-spectrum argument, implies
that every stable rational relation space is obtained from a constant vector
space.  Beukers' refinement of the Siegel--Shidlovskii theorem then converts an
algebraic value into an exponential-polynomial identity.  A
constant-coefficient annihilator and an edgewise Cauchy-transform argument
force the original exponential period to vanish identically.
\end{abstract}

\medskip
\noindent\textit{2020 Mathematics Subject Classification.}
Primary 11J81; Secondary 14F10, 34M35.

\noindent\textit{Keywords.}
$E$-functions, Siegel--Shidlovskii theorem, exponential integrals,
polynomial moments, transcendence, linear independence.

\tableofcontents

\section{Introduction and main results}\label{sec:introduction}

For a polynomial phase $q\in\Qbar[x]$ and algebraic endpoints $A,B$, consider
\begin{equation}\label{eq:basic-integral}
  \int_A^B e^{\xi q(x)}\,dx,
  \qquad \xi\in\Qbar^\times.
\end{equation}
When $q$ is linear, the integral is an algebraic multiple of a difference of
exponentials, and its transcendence follows from the
Lindemann--Weierstrass theorem.  For higher-degree phases, the integral is a
special value of an $E$-function, but direct use of the
Siegel--Shidlovskii theory is obstructed by possible functional relations among
the auxiliary moment functions needed to form a first-order differential
system.  The purpose of this paper is to resolve that obstruction and to
identify all algebraic linear relations among finite families of the values
\eqref{eq:basic-integral}.

The arithmetic input is Beukers' complete lifting theorem for relations among
values of $E$-functions at nonsingular nonzero algebraic points
\cite{Beukers}; see also the earlier linear-independence theorem of
Nesterenko and Shidlovskii \cite{NesterenkoShidlovskii}.  The general problem
of exceptional algebraic values of $E$-functions was made effective by
Adamczewski and Rivoal \cite{AdamczewskiRivoal} and developed algorithmically
through minimal differential equations by Bostan, Rivoal, and Salvy
\cite{BostanRivoalSalvy}.  Delaygue recently proved a
Lindemann--Weierstrass theorem for values of $E$-functions under a separation
condition on the finite singularities of their associated $G$-functions
\cite{Delaygue}.  The result here is complementary: the algebraic parameter
is fixed, the boundary zero-chain varies, and the full relation space is
identified without a singularity-separation hypothesis.

The structural input is a complete-reducibility theorem for the homogeneous
moment system associated with the finite map $q:\Aone_x\to\Aone_t$.  We
isolate that theorem from the arithmetic argument and formulate everything
that follows in terms of matrices, stable rational subspaces, and rational
projectors.  Its standard proof is a short application of the decomposition
theorem and Fourier transform for holonomic modules.  The final analytic step
is a Cauchy-transform rigidity argument in the spirit of the polynomial
moment problem and of Cauchy-type integrals of algebraic functions
\cite{PakovichMuzychuk,PakovichRoytvarfYomdin}.

We first formulate the results for algebraic zero-chains.  This both clarifies
the exact linear-independence statement and records the unavoidable functional
relations that can occur.

\begin{definition}[Algebraic zero-chains and exponential periods]
\label{def:zero-chain}
A finite algebraic zero-chain of degree zero is an expression
\begin{equation}\label{eq:zero-chain}
  \Delta=\sum_{\alpha\in S}d_\alpha[\alpha],
  \qquad
  S\subset\Qbar\text{ finite},\quad
  d_\alpha\in\Qbar,\quad
  \sum_{\alpha\in S}d_\alpha=0.
\end{equation}
Choose a finite piecewise $C^1$ one-chain $\Gamma$ in $\C$, with complex
coefficients, such that $\partial\Gamma=\Delta$, and define
\begin{equation}\label{eq:I-q-Delta}
  I_{q,\Delta}(z)=\int_\Gamma e^{zq(x)}\,dx.
\end{equation}
The value is independent of $\Gamma$: the one-form $e^{zq(x)}dx$ is entire in
$x$, and every closed one-chain in $\C$ has zero integral.  Thus
$I_{q,\Delta}$ depends only on $q$ and $\Delta$ and is an entire function of
$z$.
\end{definition}

The notation in \eqref{eq:I-q-Delta} may be read as integration ``over the
boundary data'' $\Delta$.  For $\Delta=[B]-[A]$, it is simply
\[
  I_{q,[B]-[A]}(z)=\int_A^B e^{zq(x)}\,dx.
\]

\begin{theorem}[Algebraic-value rigidity]\label{thm:algebraic-value-rigidity}
Let $q\in\Qbar[x]$ be nonconstant, let $\Delta$ be a finite algebraic
zero-chain of degree zero, and let $\xi\in\Qbar^\times$.  Then
\begin{equation}\label{eq:dichotomy}
  I_{q,\Delta}(\xi)\in\Qbar
  \quad\Longleftrightarrow\quad
  I_{q,\Delta}(z)\equiv0.
\end{equation}
In the equivalent cases, $I_{q,\Delta}(\xi)=0$.  Equivalently, if
$I_{q,\Delta}$ is not the zero function, then $I_{q,\Delta}(\xi)$ is
transcendental for every $\xi\in\Qbar^\times$.
\end{theorem}

\begin{corollary}[Algebraic exponential-integral theorem]
\label{cor:single-integral}
Let $q\in\Qbar[x]$ be nonconstant, let $A,B\in\Qbar$ with $A\ne B$, and let
$\xi\in\Qbar^\times$.  Then
\[
  \int_A^B e^{\xi q(x)}\,dx
\]
is transcendental.  In particular, it is nonzero.
\end{corollary}

Indeed, $I_{q,[B]-[A]}(0)=B-A\ne0$, so the corresponding entire function is
not identically zero.

The next theorem is the exact linear-independence consequence of
\cref{thm:algebraic-value-rigidity}.

\begin{theorem}[Transfer of linear relations]\label{thm:relation-transfer}
Let $q\in\Qbar[x]$ be nonconstant, let
$\Delta_1,\ldots,\Delta_r$ be finite algebraic zero-chains of degree zero, put
$I_j(z)=I_{q,\Delta_j}(z)$, and let $\xi\in\Qbar^\times$.  For
$a_0,a_1,\ldots,a_r\in\Qbar$, one has
\begin{equation}\label{eq:relation-transfer}
  a_0+\sum_{j=1}^r a_jI_j(\xi)=0
  \quad\Longleftrightarrow\quad
  a_0=0\ \text{ and }\
  \sum_{j=1}^r a_jI_j(z)\equiv0.
\end{equation}
Consequently,
\begin{equation}\label{eq:dimension-transfer}
  \dim_{\Qbar}\Span_{\Qbar}\{1,I_1(\xi),\ldots,I_r(\xi)\}
  =1+\dim_{\Qbar}\Span_{\Qbar}\{I_1,\ldots,I_r\},
\end{equation}
where the span on the right is taken in the vector space of entire functions.
\end{theorem}

Functional relations genuinely occur and therefore cannot be omitted from the
statement.  For example, if $q$ is even, then
\[
  \int_a^b e^{zq(x)}\,dx
  +\int_{-a}^{-b}e^{zq(x)}\,dx=0
\]
identically in $z$.  The following concrete criterion excludes this kind of
collision.

\begin{corollary}[Linear independence over separated regular values]
\label{cor:regular-value-independence}
Let $a_0,a_1,\ldots,a_r\in\Qbar$ be such that
$q(a_0),q(a_1),\ldots,q(a_r)$ are pairwise distinct regular values of $q$.
Then, for every $\xi\in\Qbar^\times$, the numbers
\begin{equation}\label{eq:regular-value-family}
  1,\quad
  \int_{a_0}^{a_1}e^{\xi q(x)}\,dx,\quad\ldots,\quad
  \int_{a_0}^{a_r}e^{\xi q(x)}\,dx
\end{equation}
are linearly independent over $\Qbar$.
\end{corollary}

Here a regular value means a point of the target outside the finite set of
critical values of $q$.  We prove a more precise edge-density criterion for
functional vanishing in \cref{sec:phase-rigidity};
\cref{cor:regular-value-independence} follows from it in
\cref{sec:main-proofs}.

The proof of \cref{thm:algebraic-value-rigidity} has four stages.  First, the
moments
\[
  U_j(z)=\int_\Delta x^j e^{zq(x)}\,dx,
  \qquad 0\le j\le\deg q-2,
\]
form an explicit inhomogeneous first-order system and are $E$-functions.
Second, complete reducibility and a rational-matrix rigidity theorem give a
constant-splitting result for every stable relation space of that system.
Third, Beukers' theorem and constant splitting show that an algebraic value of
$U_0$ forces $U_0$ to be an exponential polynomial.  Finally, a
constant-coefficient annihilator produces a vanishing polynomial moment
sequence, and an edgewise Cauchy-transform argument forces $U_0$ itself to
vanish identically.

\section{\texorpdfstring{$E$}{E}-functions and value-relation lifting}\label{sec:E-functions}

We use the definition adopted by Beukers \cite[Section~1]{Beukers}.

\begin{definition}[$E$-function]\label{def:E-function}
An entire function
\[
  F(z)=\sum_{m\ge0}a_m\frac{z^m}{m!}
\]
is an $E$-function if:
\begin{enumerate}[label=\textup{(\roman*)},leftmargin=2.2em]
\item $a_m\in\Qbar$ for every $m$;
\item the logarithmic height $h(a_0,\ldots,a_m)$ is $O(m)$;
\item $F$ satisfies a nonzero linear differential equation with coefficients
in $\Qbar[z]$.
\end{enumerate}
\end{definition}

The decisive arithmetic input is the following specialization theorem.

\begin{theorem}[Beukers]\label{thm:Beukers}
Let $F_1,\ldots,F_N$ be $E$-functions forming a vector solution of
\[
  \mathbf F'=A(z)\mathbf F,
  \qquad A(z)\in M_N(\Qbar(z)),
\]
and let $T(z)\in\Qbar[z]$ be a common denominator of the entries of $A(z)$.
Let $\xi\in\Qbar$ satisfy $\xi T(\xi)\ne0$.  If
$P\in\Qbar[X_1,\ldots,X_N]$ is homogeneous and
\[
  P(F_1(\xi),\ldots,F_N(\xi))=0,
\]
then there exists
$Q\in\Qbar[z,X_1,\ldots,X_N]$, homogeneous in $X_1,\ldots,X_N$ and of the
same degree as $P$, such that
\[
  Q(z,F_1(z),\ldots,F_N(z))\equiv0,
  \qquad
  Q(\xi,X_1,\ldots,X_N)=P(X_1,\ldots,X_N).
\]
In particular, a prescribed linear relation at $z=\xi$ lifts to a linear
functional relation whose coefficient polynomials specialize coefficientwise
to the prescribed relation.
\end{theorem}

This combines \cite[Theorems~1.3 and 3.2]{Beukers}.  Adjoining the constant
$E$-function $1$ converts affine relations into homogeneous ones.

We shall also use two standard independence statements.

\begin{lemma}[Distinct exponential functions]\label{lem:distinct-exponentials}
If $\lambda_1,\ldots,\lambda_s\in\C$ are distinct, then
$e^{\lambda_1z},\ldots,e^{\lambda_sz}$ are linearly independent over
$\C(z)$.
\end{lemma}

\begin{proof}
Suppose a nontrivial relation with the fewest nonzero terms exists.
Differentiate it and eliminate one term.  Minimality forces the logarithmic
derivative of a nonzero rational function to be a nonzero constant.  This is
impossible, since the Laurent expansion at infinity of a logarithmic derivative
has zero constant term.
\end{proof}

\begin{theorem}[Lindemann--Weierstrass, linear form]
\label{thm:Lindemann-Weierstrass}
If $\beta_1,\ldots,\beta_s\in\Qbar$ are distinct, then
$e^{\beta_1},\ldots,e^{\beta_s}$ are linearly independent over $\Qbar$.
\end{theorem}

We refer to \cite[Chapter~1]{Baker} for the classical theorem.

\section{The moment differential system}\label{sec:moment-system}

Throughout \cref{sec:moment-system,sec:moment-E-functions,%
sec:constant-splitting,sec:value-to-exponential}, let
$q\in\Qbar[x]$ have degree $n\ge2$.  No normalization of $q$ is needed.
For $0\le j\le n-2$, perform Euclidean division
\begin{equation}\label{eq:division}
  q(x)x^j=h_j(x)q'(x)+r_j(x),
  \qquad \deg r_j\le n-2.
\end{equation}
Let $\mathsf C$ and $\mathsf D$ be the matrices whose $j$th columns are the
coefficient vectors of $r_j$ and $-h_j'$ in the basis
$1,x,\ldots,x^{n-2}$.

\begin{lemma}[The two division matrices]\label{lem:CD}
The following assertions hold.
\begin{enumerate}[label=\textup{(\roman*)},leftmargin=2.2em]
\item $\mathsf C$ is multiplication by $q$ on the Jacobian algebra
$\Qbar[x]/(q')$;
\item $\mathsf D$ is triangular with diagonal
\[
  -\frac1n,-\frac2n,\ldots,-\frac{n-1}{n}.
\]
\end{enumerate}
\end{lemma}

\begin{proof}
The remainder $r_j$ is precisely the representative of $qx^j$ modulo $q'$
of degree at most $n-2$, which proves (i).  If the leading coefficient of
$q$ is $a_n$, then comparison of leading terms in \eqref{eq:division} gives
\[
  h_j(x)=\frac{x^{j+1}}n+
  \text{a polynomial of degree at most }j.
\]
Consequently
\[
  -h_j'(x)=-\frac{j+1}{n}x^j+
  \text{terms of degree at most }j-1.
\]
Thus $\mathsf D$ is triangular with the stated diagonal.
\end{proof}

Fix a finite algebraic zero-chain
\begin{equation}\label{eq:Delta-fixed}
  \Delta=\sum_{\alpha\in S}d_\alpha[\alpha]
\end{equation}
of degree zero.  Choose a one-chain $\Gamma$ with
$\partial\Gamma=\Delta$, and define
\begin{equation}\label{eq:Uj}
  U_j(z)=\int_\Gamma x^j e^{zq(x)}\,dx,
  \qquad
  \mathbf U=(U_0,\ldots,U_{n-2})^{\mathsf T}.
\end{equation}
Define
\begin{equation}\label{eq:b-x}
  \mathbf b(x)=n\bigl(h_0(x),\ldots,h_{n-2}(x)\bigr)^{\mathsf T}.
\end{equation}
Let
\begin{equation}\label{eq:Lambda-Delta}
  \Lambda_\Delta=
  \{q(\alpha):\alpha\in S,\ d_\alpha\ne0\}
\end{equation}
with repetitions removed, and for $\lambda\in\Lambda_\Delta$ put
\begin{equation}\label{eq:b-lambda}
  \mathbf b_\lambda
  =\sum_{\substack{\alpha\in S\\q(\alpha)=\lambda}}
    d_\alpha\mathbf b(\alpha).
\end{equation}

\begin{proposition}[Moment differential system]\label{prop:moment-system}
Put $C=\mathsf C^{\mathsf T}$ and $D=\mathsf D^{\mathsf T}$.  Then
\begin{equation}\label{eq:moment-system}
  \mathbf U'
  =\left(C+\frac Dz\right)\mathbf U
   +\frac1{nz}\sum_{\lambda\in\Lambda_\Delta}
      \mathbf b_\lambda e^{\lambda z}.
\end{equation}
Moreover,
\begin{equation}\label{eq:D-spectrum}
  \Spec(D)=
  \left\{-\frac1n,-\frac2n,\ldots,-\frac{n-1}{n}\right\}.
\end{equation}
\end{proposition}

\begin{proof}
Using \eqref{eq:division} and integration by parts along $\Gamma$, we obtain,
for $z\ne0$,
\begin{align*}
  U_j'
  &=\int_\Gamma q(x)x^j e^{zq(x)}\,dx\\
  &=\int_\Gamma\left(r_j(x)-\frac1z h_j'(x)\right)e^{zq(x)}\,dx
    +\frac1z\sum_{\alpha\in S}
      d_\alpha h_j(\alpha)e^{zq(\alpha)}.
\end{align*}
Writing these identities simultaneously gives \eqref{eq:moment-system}.
The spectral assertion follows from \cref{lem:CD}.
\end{proof}

\section{The moment functions are \texorpdfstring{$E$}{E}-functions}
\label{sec:moment-E-functions}

We record a height estimate that will be used to verify the arithmetic growth
condition in \cref{def:E-function}.

\begin{lemma}[Joint height from one common denominator]
\label{lem:joint-height}
Let $K_0$ be a number field, let $\alpha_0,\ldots,\alpha_m\in K_0$, and let
$\delta\in\mathbb Z_{\ge1}$ satisfy
$\delta\alpha_k\in\mathcal O_{K_0}$ for every $k$.  Then
\[
  h(\alpha_0,\ldots,\alpha_m)
  \le \log\delta
  +\log\max\left(
     1,\max_{0\le k\le m}\max_{\sigma:K_0\hookrightarrow\C}
     |\sigma(\alpha_k)|
   \right).
\]
Here the left side is the absolute logarithmic Weil height of
$(1:\alpha_0:\cdots:\alpha_m)$.
\end{lemma}

\begin{proof}
With the standard absolute values and local degrees
$n_v=[(K_0)_v:\Q_v]$,
\[
  h(\alpha_0,\ldots,\alpha_m)
  =\frac1{[K_0:\Q]}\sum_{v\in M_{K_0}}n_v
   \log\max(1,|\alpha_0|_v,\ldots,|\alpha_m|_v).
\]
Put $\beta_k=\delta\alpha_k\in\mathcal O_{K_0}$.  Projective invariance gives
\[
  (1:\alpha_0:\cdots:\alpha_m)
  =(\delta:\beta_0:\cdots:\beta_m).
\]
At every finite place all displayed coordinates have absolute value at most
$1$.  At an infinite place,
\[
  \max(|\delta|_v,|\beta_0|_v,\ldots,|\beta_m|_v)
  \le \delta\max\left(
     1,\max_{k,\sigma}|\sigma(\alpha_k)|
  \right).
\]
Summing over the infinite places proves the claim.
\end{proof}

\begin{lemma}[Moment functions are $E$-functions]\label{lem:moment-E-functions}
For $0\le j\le n-2$, the entire function $U_j$ in \eqref{eq:Uj} is an
$E$-function.  More precisely, if
\begin{equation}\label{eq:Uj-series}
  U_j(z)=\sum_{m\ge0}u_{j,m}\frac{z^m}{m!},
  \qquad
  u_{j,m}=\int_\Gamma x^j q(x)^m\,dx,
\end{equation}
then there is a constant $\kappa_j$ such that
\[
  h(u_{j,0},\ldots,u_{j,m})\le\kappa_j(m+1)
  \qquad(m\ge0),
\]
and $U_j$ satisfies a nonzero homogeneous linear differential equation over
$\Qbar[z]$ of order at most $n-1+|\Lambda_\Delta|$.
\end{lemma}

\begin{proof}
Uniform convergence on compact subsets permits termwise integration of
$e^{zq(x)}$ along $\Gamma$, proving \eqref{eq:Uj-series} and the entirety of
$U_j$.

Let
\[
  K_0=\Q\bigl(
    \{d_\alpha,\alpha:\alpha\in S\},
    \text{the coefficients of }q
  \bigr).
\]
Write
\[
  q(x)=\sum_{i=0}^n c_ix^i,
  \qquad
  q(x)^m=\sum_{k=0}^{nm}\gamma_{m,k}x^k.
\]
Since a primitive of $x^{j+k}$ is $x^{j+k+1}/(j+k+1)$, the boundary formula
for integration gives the exact identity
\begin{equation}\label{eq:u-exact-general}
  u_{j,m}
  =\sum_{k=0}^{nm}\gamma_{m,k}
    \frac{\displaystyle\sum_{\alpha\in S}
      d_\alpha\alpha^{j+k+1}}{j+k+1}.
\end{equation}
Thus $u_{j,m}\in K_0$.

We first bound all archimedean conjugates simultaneously.  Put
\[
  M=\max_{\substack{\sigma:K_0\hookrightarrow\C\\\alpha\in S}}
      \max(1,|\sigma(\alpha)|),
  \qquad
  \Theta=\max_{\sigma:K_0\hookrightarrow\C}
          \max_{|x|\le M}|q^\sigma(x)|,
\]
and
\[
  D_0=\max_{\sigma:K_0\hookrightarrow\C}
      \sum_{\alpha\in S}|\sigma(d_\alpha)|.
\]
If $P_{j,m}^\sigma(x)=\int_0^x t^j q^\sigma(t)^m\,dt$ is taken along the
straight segment, then applying $\sigma$ to \eqref{eq:u-exact-general} gives
\[
  \sigma(u_{j,m})
  =\sum_{\alpha\in S}\sigma(d_\alpha)
    P_{j,m}^\sigma(\sigma(\alpha)).
\]
Every such segment lies in $|x|\le M$, and hence
\[
  |\sigma(u_{j,m})|
  \le D_0 M^{j+1}\Theta^m.
\]
This is bounded by $C_1^{m+j+2}$ for a constant $C_1\ge1$ independent of
$m$ and $\sigma$.

Choose $d\in\mathbb Z_{\ge1}$ so that
\[
  dc_0,\ldots,dc_n,\quad
  d\alpha,\quad dd_\alpha
  \in\mathcal O_{K_0}
  \qquad(\alpha\in S).
\]
For $m\ge0$, set
\begin{equation}\label{eq:delta-m}
  \delta_m
  =d^{(n+1)m+j+2}
   \operatorname{lcm}(1,\ldots,nm+j+1).
\end{equation}
Then $d^m\gamma_{m,k}$, $dd_\alpha$, and
$d^{j+k+1}\alpha^{j+k+1}$ are algebraic integers.  Since $0\le k\le nm$,
each summand in \eqref{eq:u-exact-general}, after multiplication by
$\delta_m$, is a sum of products of algebraic integers; the spare factor
$d^{nm-k}$ supplies the uniform exponent in \eqref{eq:delta-m}.  Therefore
$\delta_mu_{j,m}$ is integral.  Moreover, $\delta_\ell\mid\delta_m$ for
$0\le\ell\le m$, so $\delta_m$ is a common denominator for
$u_{j,0},\ldots,u_{j,m}$.

By \cref{lem:joint-height},
\[
  h(u_{j,0},\ldots,u_{j,m})
  \le\log\delta_m+(m+j+2)\log C_1.
\]
The elementary Chebyshev estimate
$\log\operatorname{lcm}(1,\ldots,N)\le2N$ gives
$\log\delta_m=O(m)$, and hence the required joint-height bound.

It remains to establish a scalar differential equation.  Let
$\lambda_1,\ldots,\lambda_s$ be the distinct elements of
$\Lambda_\Delta$.  By \cref{prop:moment-system}, the $\Qbar(z)$-span of
\[
  U_0,\ldots,U_{n-2},
  e^{\lambda_1z},\ldots,e^{\lambda_sz}
\]
is stable under $\partial_z$ and has dimension at most $n-1+s$.  Therefore
$U_j$ and its first $n-1+s$ derivatives are linearly dependent over
$\Qbar(z)$.  Clearing denominators yields a nonzero operator in
$\Qbar[z]\langle\partial_z\rangle$ of order at most $n-1+s$ annihilating
$U_j$.  All three defining conditions of an $E$-function now follow.
\end{proof}

\section{Complete reducibility and constant splitting}
\label{sec:constant-splitting}

Put
\begin{equation}\label{eq:A-matrix}
  A(z)=C+\frac Dz
\end{equation}
and define, on rational row vectors,
\begin{equation}\label{eq:T-operator}
  \mathcal T(\boldsymbol\phi)
  =\boldsymbol\phi'+\boldsymbol\phi A(z).
\end{equation}
A $\C(z)$-subspace of $\C(z)^{1\times(n-1)}$ will be called stable if it is
preserved by $\mathcal T$.  We call the system completely reducible if every
stable subspace has a stable complement.

\begin{proposition}[Complete reducibility]\label{prop:complete-reducibility}
The row system \eqref{eq:T-operator} is completely reducible.
\end{proposition}

\begin{proof}
We give the standard geometric proof, but the remainder of the paper uses only
the statement of the proposition.  Let $\Sigma$ be the finite set of critical
values of $q$, put $V=\Aone_t\setminus\Sigma$, and let $\Lcal$ be the
permutation local system of the unramified covering
$q:q^{-1}(V)\to V$.  The covering is connected, so its monodromy action is
transitive.  The trace map splits
\[
  \Lcal\simeq\C_V\oplus\Lcal^0,
\]
and $\Lcal^0$ is semisimple by Maschke's theorem, since the monodromy image is
a finite subgroup of $S_n$.

The map $q:\Aone_x\to\Aone_t$ is finite, proper, and small.  Hence
\[
  Rq_*\C_{\Aone_x}[1]\simeq\operatorname{IC}(\Lcal);
\]
see \cite[Remark~2.13]{deCataldoMigliorini}.  Under the Riemann--Hilbert
correspondence and compatibility with proper direct image, this yields
\begin{equation}\label{eq:direct-image-splitting}
  q_+\Ocal_{\Aone_x}
  \simeq
  \Ocal_{\Aone_t}\oplus\Ncal^0,
\end{equation}
where $\Ncal^0$ is a direct sum of simple regular holonomic modules; see
\cite[Sections~3.4, 7.1--7.2, and 8.2.2]{HTT}.

Use a Fourier variable $\tau$ with
\begin{equation}\label{eq:Fourier-convention}
  t\longmapsto-\partial_\tau,
  \qquad
  \partial_t\longmapsto\tau.
\end{equation}
This Weyl-algebra automorphism is exact and preserves simple summands.  The
Fourier transform of $\Ocal_{\Aone_t}$ is supported at $\tau=0$ and vanishes
after localization to $\C(\tau)$.  We also record the elementary localization
point.  If $M$ is a simple Weyl module and
$S=\C[\tau]\setminus\{0\}$, then the $S$-torsion of $M$ is a Weyl submodule:
from $p(\tau)m=0$ one obtains $p(\tau)^2\partial_\tau m=0$.  Hence either
$S^{-1}M=0$, or $M$ embeds in $S^{-1}M$.  In the latter case, every nonzero
submodule of $S^{-1}M$ meets $M$ after clearing a denominator, and simplicity
forces that submodule to be all of $S^{-1}M$.  Thus the surviving generic
Fourier transform of $\Ncal^0$ is completely reducible.

For completeness, we identify it explicitly.  The graph module of $q$ is
generated by a symbol $\delta_q$ satisfying
\[
  (t-q(x))\delta_q=0,
  \qquad
  (\partial_x+q'(x)\partial_t)\delta_q=0.
\]
After \eqref{eq:Fourier-convention}, its transform satisfies
\[
  (\partial_\tau+q(x))\widehat\delta_q=0,
  \qquad
  (\partial_x+\tau q'(x))\widehat\delta_q=0.
\]
On coefficient polynomials multiplying $\widehat\delta_q$, the two relevant
operators are therefore $\partial_x-\tau q'(x)$ and
$\partial_\tau-q(x)$.  Projection along the $x$-line is computed by the
relative complex
\[
  \C(\tau)[x]
  \xrightarrow{\ \partial_x-\tau q'\ }
  \C(\tau)[x].
\]
Its kernel is zero, because for nonzero $p\in\C(\tau)[x]$ the term
$-\tau q'p$ has strictly larger $x$-degree than $p'$.  Hence the generic
transform is
\begin{equation}\label{eq:generic-transform}
  \C(\tau)[x]\Big/
  (\partial_x-\tau q'(x))\C(\tau)[x],
  \qquad
  \partial_\tau[p]=[\partial_\tau p-qp].
\end{equation}
Pulling back by $\tau=-z$ gives
\begin{equation}\label{eq:generic-z-quotient}
  \C(z)[x]\Big/
  (\partial_x+zq'(x))\C(z)[x],
  \qquad
  \partial_z[p]=[\partial_zp+qp].
\end{equation}
The classes of $1,x,\ldots,x^{n-2}$ form a basis: the relation
$[zq'p]=-[p']$ reduces every higher power, while the image of a nonzero
polynomial has degree larger by $n-1$ and therefore contains no nonzero
polynomial of degree at most $n-2$.  In this basis, \eqref{eq:division} gives
\[
  \partial_z[x^j]=[qx^j]=[r_j]-\frac1z[h_j'],
\]
so the matrix is $\mathsf C+\mathsf D/z$.  The moment matrix is its transpose,
$A(z)=C+D/z$.  Complete reducibility is preserved by transposition and by
passing to the dual row system, proving the proposition.
\end{proof}

The next result is elementary and is the only special property of the
matrices $C,D$ needed to turn complete reducibility into constant splitting.
Let
\[
  \Vcal=\C[x]_{<n-1}.
\]
For $g\in\Vcal$, divide
\begin{equation}\label{eq:g-division}
  gq=h_gq'+r_g,
  \qquad \deg r_g<n-1,
\end{equation}
and define
\begin{equation}\label{eq:CstarDstar}
  \mathsf C_*g=r_g,
  \qquad
  \mathsf D_*g=-h_g'.
\end{equation}
These operators are represented by $\mathsf C$ and $\mathsf D$ in the
monomial basis.  For a critical point $\tau$, put
\[
  m_\tau=\ord_\tau q',
  \qquad
  \lambda_\tau=q(\tau).
\]

\begin{lemma}[Critical-point basis and local spectrum]
\label{lem:critical-basis}
For every critical point $\tau$ and $1\le k\le m_\tau$, define
\begin{equation}\label{eq:g-tau-k}
  g_{\tau,k}(x)=\frac{q'(x)}{(x-\tau)^k}.
\end{equation}
Then:
\begin{enumerate}[label=\textup{(\roman*)},leftmargin=2.2em]
\item the polynomials $g_{\tau,k}$, over all $\tau$ and $k$, form a basis of
$\Vcal$;
\item $\mathsf C_*$ is semisimple and
\[
  \mathsf C_*g_{\tau,k}=\lambda_\tau g_{\tau,k};
\]
\item put
\[
  E_\tau=\Span_\C\{g_{\tau,1},\ldots,g_{\tau,m_\tau}\}.
\]
If $E_\lambda=\ker(\mathsf C_*-\lambda I)$ and $\pi_\lambda$ is the spectral
projection, then
\[
  E_\lambda=\bigoplus_{\tau:\lambda_\tau=\lambda}E_\tau,
  \qquad
  \overline D_\lambda
  =\pi_\lambda\mathsf D_*|_{E_\lambda}
\]
preserves every summand $E_\tau$.  In the ordered basis
$g_{\tau,1},\ldots,g_{\tau,m_\tau}$, its restriction to $E_\tau$ is
triangular with diagonal
\begin{equation}\label{eq:local-D-spectrum}
  -\frac{m_\tau}{m_\tau+1},
  -\frac{m_\tau-1}{m_\tau+1},
  \ldots,
  -\frac1{m_\tau+1}.
\end{equation}
\end{enumerate}
\end{lemma}

\begin{proof}
Factor
\[
  q'(x)=c\prod_\tau(x-\tau)^{m_\tau}.
\]
Reduction modulo $q'$ identifies $\Vcal$ with $\C[x]/(q')$, and the Chinese
remainder theorem gives
\[
  \C[x]/(q')
  \simeq
  \bigoplus_\tau\C[x]/(x-\tau)^{m_\tau}.
\]
For fixed $\tau$, the images of $g_{\tau,k}$ vanish in every other local
factor, while in the $\tau$-factor, with $y=x-\tau$, they are
\[
  g_{\tau,k}=u_\tau(y)y^{m_\tau-k},
  \qquad u_\tau(0)\ne0.
\]
Hence $g_{\tau,1},\ldots,g_{\tau,m_\tau}$ form a basis of the $\tau$-factor.
Since $\sum_\tau m_\tau=n-1$, all the displayed vectors form a basis of
$\Vcal$.

Since $q-\lambda_\tau$ vanishes at $\tau$ to order $m_\tau+1$, the polynomial
\[
  h_{\tau,k}(x)=\frac{q(x)-\lambda_\tau}{(x-\tau)^k}
\]
satisfies
\[
  qg_{\tau,k}=h_{\tau,k}q'+\lambda_\tau g_{\tau,k}.
\]
Thus $\mathsf C_*g_{\tau,k}=\lambda_\tau g_{\tau,k}$, which proves (ii) and
the displayed decomposition of $E_\lambda$.

Write $m=m_\tau$, $y=x-\tau$, and
\[
  q(x)-\lambda_\tau=ay^{m+1}+O(y^{m+2}),
  \qquad a\ne0.
\]
Then
\[
  g_{\tau,k}=a(m+1)y^{m-k}+O(y^{m-k+1})
\]
and
\[
  -h_{\tau,k}'
  =-a(m+1-k)y^{m-k}+O(y^{m-k+1}).
\]
In the ordered local basis, the error term involves only
$g_{\tau,j}$ with $j<k$.  Hence the compressed local block is triangular and
its $k$th diagonal entry is $-(m+1-k)/(m+1)$.

If $\sigma\ne\tau$ and $\lambda_\sigma=\lambda_\tau$, then
$h_{\tau,k}=(q-\lambda_\tau)/(x-\tau)^k$ vanishes at $\sigma$ to order
$m_\sigma+1$.  Hence $h_{\tau,k}'$ vanishes modulo
$(x-\sigma)^{m_\sigma}$.  This excludes cross terms between distinct critical
points over the same critical value and proves (iii).
\end{proof}

\begin{theorem}[Rational symmetry rigidity]\label{thm:end-rigidity}
Let $P(z)\in M_{n-1}(\C(z))$ satisfy
\begin{equation}\label{eq:endomorphism-equation}
  P'=[A(z),P].
\end{equation}
Then $P$ is constant and commutes with both $C$ and $D$.
\end{theorem}

\begin{proof}
First exclude finite poles.  If $P$ had a pole of order $m\ge1$ at
$z_0\ne0$, then $P'$ would have pole order $m+1$, whereas the right side of
\eqref{eq:endomorphism-equation} would have order at most $m$.  Thus no such
pole occurs.

At $z=0$, suppose
\[
  P=z^{-m}P_{-m}+O(z^{-m+1}),
  \qquad m\ge1,
  \quad P_{-m}\ne0.
\]
The coefficient of $z^{-m-1}$ is
\[
  -mP_{-m}=[D,P_{-m}].
\]
By \eqref{eq:D-spectrum}, every eigenvalue of
$\operatorname{ad}D:X\mapsto[D,X]$ lies strictly between $-1$ and $1$, so
$-m$ cannot occur.  Hence $P$ has no finite pole and is a polynomial matrix.

Transpose \eqref{eq:endomorphism-equation} and put $Q=P^{\mathsf T}$.  Then
\begin{equation}\label{eq:row-endomorphism}
  Q'=[Q,\mathsf C_*+\mathsf D_*/z].
\end{equation}
Suppose
\[
  Q(z)=Q_dz^d+Q_{d-1}z^{d-1}+\cdots,
  \qquad d\ge1,
  \quad Q_d\ne0.
\]
Multiplying \eqref{eq:row-endomorphism} by $z$ and comparing the coefficients
of $z^{d+1}$ and $z^d$ gives
\begin{align}
  [Q_d,\mathsf C_*]&=0,\label{eq:leading-commute}\\
  dQ_d&=[Q_{d-1},\mathsf C_*]+[Q_d,\mathsf D_*].
  \label{eq:next-leading}
\end{align}

For each eigenvalue $\lambda$ of $\mathsf C_*$, the spectral projection
$\pi_\lambda$ is a polynomial in $\mathsf C_*$.  Hence $Q_d$ commutes with
$\pi_\lambda$ and is block diagonal.  For a block
$Q_{d,\lambda}=\pi_\lambda Q_d\pi_\lambda\ne0$, compressing
\eqref{eq:next-leading} to $E_\lambda$ gives
\begin{equation}\label{eq:ad-local-D}
  dQ_{d,\lambda}
  =[Q_{d,\lambda},\overline D_\lambda].
\end{equation}
After triangularizing $\overline D_\lambda$, the commutator on the right is
triangular on matrix units, with diagonal entries equal to pairwise
differences of eigenvalues of $\overline D_\lambda$.  By
\eqref{eq:local-D-spectrum}, all these differences lie strictly between
$-1$ and $1$.  Equation \eqref{eq:ad-local-D} would make the positive integer
$d$ an eigenvalue, a contradiction.  Thus $d=0$.

Therefore $P$ is constant.  Substitution into
\eqref{eq:endomorphism-equation} gives
$[C,P]+z^{-1}[D,P]=0$, hence $[C,P]=[D,P]=0$.
\end{proof}

\begin{theorem}[Constant splitting]\label{thm:constant-splitting}
Let $\Rcal\subseteq\C(z)^{1\times(n-1)}$ be stable under $\mathcal T$.  Then
there is a constant subspace $W\subseteq\C^{1\times(n-1)}$, invariant under
right multiplication by $C$ and by $D$, such that
\begin{equation}\label{eq:constant-splitting}
  \Rcal=W\otimes_\C\C(z).
\end{equation}
\end{theorem}

\begin{proof}
By \cref{prop:complete-reducibility}, choose a stable complement of $\Rcal$.
Let right multiplication by $P(z)$ be the projector onto $\Rcal$ along that
complement.  Since both summands are stable, the projector commutes with
$\mathcal T$.  For an arbitrary row $\boldsymbol\phi$, this identity is
\[
  \mathcal T(\boldsymbol\phi P)=\mathcal T(\boldsymbol\phi)P,
\]
which is equivalent to
\[
  P'=[A(z),P].
\]
By \cref{thm:end-rigidity}, $P$ is constant and commutes with $C,D$.  Taking
$W$ to be the constant row space of $P$ gives
\eqref{eq:constant-splitting}; commutation with $C,D$ gives the stated
invariance.
\end{proof}

\section{From an algebraic value to an exponential polynomial}
\label{sec:value-to-exponential}

Let
\begin{equation}\label{eq:Lambda-all}
  \Lambda=\{0\}\cup\Lambda_\Delta
\end{equation}
with repetitions removed, and put
\begin{equation}\label{eq:E-space}
  \Ecal=
  \Span_{\C(z)}\{e^{\lambda z}:\lambda\in\Lambda\}.
\end{equation}
If $0\notin\Lambda_\Delta$, its forcing vector is understood to be zero:
$\mathbf b_0=0$.

The functions
\[
  U_0,\ldots,U_{n-2},
  \{e^{\lambda z}:\lambda\in\Lambda\}
\]
form a vector of $E$-functions.  Adjoining the exponential coordinates to
\eqref{eq:moment-system} gives a homogeneous block system over $\Qbar(z)$
with common denominator $z$.  The coordinate $e^{0z}=1$ homogenizes an
algebraic value relation.

Define the rational relation space
\begin{equation}\label{eq:relation-space}
  \Rcal=
  \left\{
    \boldsymbol\phi\in\C(z)^{1\times(n-1)}:
    \boldsymbol\phi\mathbf U\in\Ecal
  \right\}.
\end{equation}
It is stable under $\mathcal T$.  Indeed, if
\[
  \mathbf f(z)=\frac1{nz}
  \sum_{\lambda\in\Lambda}\mathbf b_\lambda e^{\lambda z},
\]
then \eqref{eq:moment-system} gives
\[
  (\boldsymbol\phi\mathbf U)'
  =\mathcal T(\boldsymbol\phi)\mathbf U
   +\boldsymbol\phi\mathbf f.
\]
Both the left side and the final term lie in $\Ecal$, so
$\mathcal T(\boldsymbol\phi)\in\Rcal$.

\begin{proposition}[An algebraic value forces an exponential polynomial]
\label{prop:value-to-exp-poly}
Let $\xi\in\Qbar^\times$.  If $U_0(\xi)\in\Qbar$, then
\begin{equation}\label{eq:exp-polynomial-U0}
  U_0(z)=\sum_{\lambda\in\Lambda}p_\lambda(z)e^{\lambda z}
\end{equation}
for polynomials $p_\lambda\in\C[z]$.
\end{proposition}

\begin{proof}
Let $a=U_0(\xi)\in\Qbar$, order
\[
  \Lambda=\{0,\lambda_2,\ldots,\lambda_s\},
\]
and form
\[
  \mathbf F(z)
  =\bigl(
    U_0(z),\ldots,U_{n-2}(z),
    1,e^{\lambda_2z},\ldots,e^{\lambda_sz}
  \bigr)^{\mathsf T}.
\]
Put
\[
  B_\Lambda=\frac1n
  \begin{pmatrix}
    \mathbf b_0&\mathbf b_{\lambda_2}&\cdots&\mathbf b_{\lambda_s}
  \end{pmatrix},
  \qquad
  \Delta_\Lambda=\diag(0,\lambda_2,\ldots,\lambda_s).
\]
Then
\begin{equation}\label{eq:full-block-system}
  \mathbf F'
  =\begin{pmatrix}
      C+D/z&B_\Lambda/z\\
      0&\Delta_\Lambda
    \end{pmatrix}\mathbf F.
\end{equation}
The relation $U_0(\xi)-a=0$ is homogeneous in the coordinates of
$\mathbf F(\xi)$.  More explicitly, let
\[
  \boldsymbol\alpha
  =\bigl(1,0,\ldots,0;\,-a,0,\ldots,0\bigr).
\]
Then $\boldsymbol\alpha\mathbf F(\xi)=0$.  The matrix in
\eqref{eq:full-block-system} has common denominator $z$, and
$\xi\ne0$.  By \cref{thm:Beukers}, there is a polynomial coefficient row
\[
  R(z)\in\Qbar[z]^{1\times((n-1)+s)}
\]
such that
\[
  R(z)\mathbf F(z)=0,
  \qquad
  R(\xi)=\boldsymbol\alpha.
\]
Write
\[
  R(z)=\bigl(
    \boldsymbol\phi(z);
    c_0(z),c_{\lambda_2}(z),\ldots,c_{\lambda_s}(z)
  \bigr).
\]
Then
\begin{equation}\label{eq:lifted-relation}
  \boldsymbol\phi(z)\mathbf U(z)
  +c_0(z)+\sum_{j=2}^s c_{\lambda_j}(z)e^{\lambda_jz}=0,
  \qquad
  \boldsymbol\phi(\xi)=\mathbf e_0:=(1,0,\ldots,0).
\end{equation}
Thus $\boldsymbol\phi\in\Rcal$.

By \cref{thm:constant-splitting},
$\Rcal=W\otimes_\C\C(z)$ for a constant subspace $W$ invariant under $C,D$.
Equivalently, $\Rcal$ is the image of a constant projector.  Since
$\boldsymbol\phi(\xi)=\mathbf e_0$, we obtain $\mathbf e_0\in W$ and hence
\[
  U_0=\mathbf e_0\mathbf U\in\Ecal.
\]

It remains to show that the rational exponential coefficients are
polynomials.  Choose a constant full-row-rank matrix $S$ whose rows form a
basis of $W$, and put $\mathbf V=S\mathbf U$.  Invariance of $W$ gives unique
constant matrices $C_W,D_W$ such that
\[
  SC=C_WS,
  \qquad
  SD=D_WS.
\]
With $\mathbf b_{W,\lambda}=S\mathbf b_\lambda$, the system
\eqref{eq:moment-system} restricts to
\begin{equation}\label{eq:restricted-system}
  \mathbf V'
  =\left(C_W+\frac{D_W}{z}\right)\mathbf V
   +\frac1{nz}\sum_{\lambda\in\Lambda}
     \mathbf b_{W,\lambda}e^{\lambda z}.
\end{equation}
Every component of $\mathbf V$ lies in $\Ecal$.  By
\cref{lem:distinct-exponentials}, there are unique rational vectors
$\mathbf r_\lambda(z)\in\C(z)^{\dim W}$ such that
\begin{equation}\label{eq:V-expansion}
  \mathbf V(z)=\sum_{\lambda\in\Lambda}
  \mathbf r_\lambda(z)e^{\lambda z}.
\end{equation}
Equating exponential coefficients in \eqref{eq:restricted-system} gives
\begin{equation}\label{eq:restricted-block}
  \mathbf r_\lambda'
  =(C_W-\lambda I)\mathbf r_\lambda
   +\frac{D_W}{z}\mathbf r_\lambda
   +\frac1{nz}\mathbf b_{W,\lambda}.
\end{equation}

A pole of $\mathbf r_\lambda$ at $z_0\ne0$ is impossible, because
differentiation raises its order by one while the right side of
\eqref{eq:restricted-block} does not.  At $z=0$, suppose
\[
  \mathbf r_\lambda=z^{-m}v+O(z^{-m+1}),
  \qquad m\ge1,
  \quad v\ne0.
\]
The coefficient of $z^{-m-1}$ gives
\[
  -mv=D_Wv.
\]
Every eigenvalue of $D_W$ is an eigenvalue of $D$: if
$u^{\mathsf T}D_W=\mu u^{\mathsf T}$, then
$u^{\mathsf T}S\ne0$ and
\[
  (u^{\mathsf T}S)D
  =u^{\mathsf T}D_WS
  =\mu(u^{\mathsf T}S).
\]
By \eqref{eq:D-spectrum}, $D_W$ has no negative integral eigenvalue.  Hence
$\mathbf r_\lambda$ has no finite pole and belongs to
$\C[z]^{\dim W}$.

Finally choose a constant row $a_0$ with $\mathbf e_0=a_0S$.  Then
\[
  U_0=\mathbf e_0\mathbf U=a_0\mathbf V
  =\sum_{\lambda\in\Lambda}
    \bigl(a_0\mathbf r_\lambda(z)\bigr)e^{\lambda z},
\]
which is \eqref{eq:exp-polynomial-U0} with
$p_\lambda=a_0\mathbf r_\lambda\in\C[z]$.
\end{proof}

\begin{remark}\label{rem:entire-rational-coefficients}
Entirety alone does not imply polynomial exponential coefficients:
\[
  \frac{e^z-1}{z}=\frac1z e^z-\frac1z
\]
is entire although both displayed coefficients have a pole.  This is the
integral $\int_0^1e^{zx}\,dx$ for a linear phase.  For $n\ge2$, the residue
nonresonance in \eqref{eq:D-spectrum} is therefore load-bearing.
\end{remark}

\section{Constant-coefficient annihilation}\label{sec:annihilation}

Assume \eqref{eq:exp-polynomial-U0}.  Choose an integer $N$ larger than the
degrees of all nonzero polynomials $p_\lambda$; with the convention
$\deg0=-\infty$, choose
\[
  N>\max_{\lambda\in\Lambda}\deg p_\lambda.
\]
Define the constant-coefficient differential operator
\begin{equation}\label{eq:L-operator}
  L=\prod_{\lambda\in\Lambda}(\partial_z-\lambda)^N.
\end{equation}
Then $LU_0=0$.  Since differentiation under the integral sign multiplies the
integrand by $q(x)$,
\begin{equation}\label{eq:annihilated-integral}
  0=LU_0(z)
  =\int_\Gamma H(q(x))e^{zq(x)}\,dx,
\end{equation}
where
\begin{equation}\label{eq:H-def}
  H(t)=\prod_{\lambda\in\Lambda}(t-\lambda)^N
\end{equation}
is nonzero.  Expanding at $z=0$ gives
\begin{equation}\label{eq:phase-moments}
  \int_\Gamma q(x)^mH(q(x))\,dx=0
  \qquad(m\ge0).
\end{equation}
The factor $t^N$ in $H$ arises from the constant coordinate used to
homogenize the value relation; it need not have any geometric significance for
$q$ or $\Delta$.

\section{Cauchy-transform rigidity for phase moments}
\label{sec:phase-rigidity}

We now prove that the moment identities \eqref{eq:phase-moments} force the
underlying exponential period to vanish identically.  The argument also gives
an exact geometric criterion for functional relations.

\begin{lemma}[Constellation representation]\label{lem:constellation}
Let $q\in\C[x]$ be nonconstant and let
$\Delta=\sum_{\alpha\in S}d_\alpha[\alpha]$ be a finite complex zero-chain of
degree zero.  There exist:
\begin{enumerate}[label=\textup{(\roman*)},leftmargin=2.2em]
\item a finite embedded tree $T$ in the $t$-plane containing every finite
critical value of $q$ and every value $q(\alpha)$, $\alpha\in S$;
\item a finite one-chain $\widetilde\Gamma$, with complex coefficients, in the
graph $q^{-1}(T)$ with
$\partial\widetilde\Gamma=\Delta$;
\end{enumerate}
such that, on every open edge $e$ of $T$, the chain
$\widetilde\Gamma$ is a finite complex linear combination of analytic inverse
branches of $q$ over $e$.
\end{lemma}

\begin{proof}
Choose a regular center $c$ and pairwise disjoint analytic arcs from $c$ to
the distinct finite critical values and to the distinct endpoint values
$q(\alpha)$, omitting repeated arms.  Their union is an embedded star and
hence a finite tree $T$.

Over $c$, the preimage consists of the $n=\deg q$ sheets of the covering.
The lifts of the critical-value arms join these sheets according to the branch
monodromy generators.  Since the polynomial covering is connected, its
monodromy action is transitive, so the finite graph $q^{-1}(T)$ is connected.
This is the standard polynomial constellation construction; compare
\cite[Section~2.1]{PakovichMuzychuk}.

Fix a vertex $x_*$ of $q^{-1}(T)$.  For every $\alpha\in S$, choose an edge
path $P_\alpha$ in $q^{-1}(T)$ from $x_*$ to $\alpha$, and put
\[
  \widetilde\Gamma=\sum_{\alpha\in S}d_\alpha P_\alpha.
\]
Because $\sum_\alpha d_\alpha=0$, its boundary is
\[
  \partial\widetilde\Gamma
  =\sum_\alpha d_\alpha[\alpha]
   -\left(\sum_\alpha d_\alpha\right)[x_*]
  =\Delta.
\]
On the interior of every edge of $T$, the map $q$ is unramified, so the stated
inverse-branch decomposition follows.
\end{proof}

Fix such a tree and graph chain.  Orient every open edge $e$ of $T$.  If
$n=\deg q$, let
\[
  x_{e,1},\ldots,x_{e,n}:e\longrightarrow\C,
  \qquad q(x_{e,i}(t))=t,
\]
be the analytic inverse branches.  Let $f_{e,i}\in\C$ be the coefficient with
which $\widetilde\Gamma$ traverses the lift $x_{e,i}(e)$, relative to the
chosen orientation, and define the edge density
\begin{equation}\label{eq:edge-density}
  \psi_e(t)
  =\sum_{i=1}^n f_{e,i}x_{e,i}'(t)
  =\sum_{i=1}^n\frac{f_{e,i}}{q'(x_{e,i}(t))}.
\end{equation}
These densities are analytic on the open edges.  Their endpoint singularities
are integrable.  Indeed, if $t=c$ is an endpoint of $e$, a branch tends to a
point $\zeta\in q^{-1}(c)$ of ramification index $r$, and
\[
  q(x)-c=(x-\zeta)^ru(x),
  \qquad u(\zeta)\ne0.
\]
Along the edge,
\[
  x_{e,i}(t)-\zeta=O((t-c)^{1/r}),
  \qquad
  x_{e,i}'(t)=O(|t-c|^{1/r-1}),
\]
and $1/r-1>-1$.

\begin{theorem}[Phase-moment rigidity and edge-density criterion]
\label{thm:phase-moment-rigidity}
Let $q\in\C[x]$ be nonconstant, let $\Delta$ be a finite complex zero-chain of
degree zero, let $\Gamma$ be any one-chain with boundary $\Delta$, and let
$0\ne H\in\C[t]$.  Suppose
\begin{equation}\label{eq:general-phase-moments}
  \int_\Gamma q(x)^mH(q(x))\,dx=0
  \qquad(m\ge0).
\end{equation}
Then every edge density \eqref{eq:edge-density} of the constellation
representation in \cref{lem:constellation} vanishes identically.  Consequently,
\begin{equation}\label{eq:period-zero-conclusion}
  I_{q,\Delta}(z)\equiv0.
\end{equation}
In particular, for a fixed constellation representation,
\begin{equation}\label{eq:edge-criterion}
  I_{q,\Delta}(z)\equiv0
  \quad\Longleftrightarrow\quad
  \psi_e\equiv0\ \text{ for every open edge }e\subset T.
\end{equation}
\end{theorem}

\begin{proof}
Every differential $q(x)^mH(q(x))dx$ is polynomial and hence has a polynomial
primitive.  Its integral depends only on the boundary.  Therefore
\eqref{eq:general-phase-moments} remains valid after replacing $\Gamma$ by the
graph chain $\widetilde\Gamma$ from \cref{lem:constellation}.

For $s\in\C\setminus T$, define the Cauchy transform
\begin{equation}\label{eq:phase-Cauchy-transform}
  \mathcal C_H(s)
  =\int_{\widetilde\Gamma}
    \frac{H(q(x))}{s-q(x)}\,dx.
\end{equation}
The path is compact and piecewise analytic, and the denominator is bounded
away from zero locally uniformly on $\C\setminus T$.  Thus $\mathcal C_H$ is
holomorphic there.  Put
\[
  M=\max_{x\in\widetilde\Gamma}|q(x)|.
\]
For $|s|>M$, the geometric series
\[
  \frac1{s-q(x)}=\sum_{m\ge0}\frac{q(x)^m}{s^{m+1}}
\]
converges uniformly on $\widetilde\Gamma$.  Using
\eqref{eq:general-phase-moments}, we obtain
\[
  \mathcal C_H(s)
  =\sum_{m\ge0}\frac1{s^{m+1}}
    \int_{\widetilde\Gamma}q(x)^mH(q(x))\,dx
  =0
  \qquad(|s|>M).
\]
The complement of a finite embedded tree in the plane is connected.  Hence
the identity theorem gives
\begin{equation}\label{eq:Cauchy-zero}
  \mathcal C_H\equiv0
  \qquad\text{on }\C\setminus T.
\end{equation}

Changing variables $t=q(x)$ separately on every lifted edge gives the
absolutely convergent representation
\begin{equation}\label{eq:edge-Cauchy}
  \mathcal C_H(s)
  =\sum_e\int_e\frac{H(t)\psi_e(t)}{s-t}\,dt,
  \qquad s\in\C\setminus T.
\end{equation}
Fix an open edge $e$ and an interior point $t_0\in e$ with $H(t_0)\ne0$.
Choose a compact analytic subarc $\eta\Subset e$ containing $t_0$ in its
interior and a disk $D$ about $t_0$ such that
\[
  \overline D\cap(T\setminus\eta)=\varnothing.
\]
After shortening $\eta$, all inverse branches occurring in $\psi_e$ extend
holomorphically to a neighborhood of $\eta$.  Thus
\[
  \rho(t)=H(t)\psi_e(t)
\]
is holomorphic there.  The contribution to \eqref{eq:edge-Cauchy} from
$T\setminus\eta$ extends holomorphically across $D$, so the entire lateral jump
of $\mathcal C_H$ across $\eta$ is the jump of
\[
  J(s)=\int_\eta\frac{\rho(t)}{s-t}\,dt.
\]

Let $a,b$ be the initial and terminal endpoints of the oriented arc $\eta$.
For $s$ near $t_0$, define
\[
  K(t,s)=
  \begin{cases}
    \dfrac{\rho(t)-\rho(s)}{t-s},&t\ne s,\\[2mm]
    \rho'(s),&t=s.
  \end{cases}
\]
Then $K$ is jointly holomorphic near $\eta\times\eta$, and
\begin{equation}\label{eq:local-Cauchy-splitting}
  J(s)
  =\rho(s)\int_\eta\frac{dt}{s-t}
   -\int_\eta K(t,s)\,dt.
\end{equation}
The second term is holomorphic across $\eta$.  On either side of the arc, the
first scalar integral is a branch of
\[
  \int_\eta\frac{dt}{s-t}
  =\log(s-a)-\log(s-b).
\]
Its two lateral determinations differ by
$2\pi i\varepsilon_e$, where $\varepsilon_e\in\{\pm1\}$ depends only on the
orientation and the designation of the two sides.  Hence
\begin{equation}\label{eq:local-jump}
  (\mathcal C_H)_+(t_0)-(\mathcal C_H)_-(t_0)
  =2\pi i\varepsilon_e H(t_0)\psi_e(t_0).
\end{equation}
By \eqref{eq:Cauchy-zero}, both lateral values vanish.  Since $H(t_0)\ne0$,
we conclude that $\psi_e(t_0)=0$.  The point $t_0$ was arbitrary outside the
finite zero set of $H$, and $\psi_e$ is analytic on the open edge.  Therefore
\[
  \psi_e\equiv0
\]
for every open edge $e$.

Finally, changing variables edge by edge in the exponential integral gives
\[
  I_{q,\Delta}(z)
  =\int_{\widetilde\Gamma}e^{zq(x)}\,dx
  =\sum_e\int_e e^{zt}\psi_e(t)\,dt
  =0.
\]
This proves \eqref{eq:period-zero-conclusion}.  If
$I_{q,\Delta}\equiv0$, its Taylor coefficients at $z=0$ give
\eqref{eq:general-phase-moments} with $H=1$, so the forward implication in
\eqref{eq:edge-criterion} follows from the theorem.  The reverse implication
is the final displayed edge decomposition.
\end{proof}

\section{Proofs of the main theorems}\label{sec:main-proofs}

\begin{proof}[Proof of \cref{thm:algebraic-value-rigidity}]
The implication
$I_{q,\Delta}\equiv0\Rightarrow I_{q,\Delta}(\xi)\in\Qbar$ is immediate.
We prove the converse.

Suppose first that $q(x)=ax+b$ with $a\ne0$.  Since
$e^{zq(x)}/(az)$ is a primitive for $z\ne0$, one has
\begin{equation}\label{eq:linear-phase-divisor}
  I_{q,\Delta}(z)
  =\frac1{az}\sum_{\alpha\in S}d_\alpha e^{zq(\alpha)}.
\end{equation}
Assume that $I_{q,\Delta}(\xi)=\gamma\in\Qbar$.  Then
\begin{equation}\label{eq:linear-LW-relation}
  \sum_{\alpha\in S}d_\alpha e^{\xi q(\alpha)}
  -a\xi\gamma e^0=0.
\end{equation}
The map $q$ is injective, so the exponents $\xi q(\alpha)$ are pairwise
distinct.  After combining the coefficient of $e^0$ if one of them is zero,
\cref{thm:Lindemann-Weierstrass} shows that $d_\alpha=0$ whenever
$q(\alpha)\ne0$.  If there is an $\alpha_0$ with $q(\alpha_0)=0$, the only
remaining coefficient relation is
$d_{\alpha_0}=a\xi\gamma$; but the degree-zero condition
$\sum_\alpha d_\alpha=0$ gives $d_{\alpha_0}=0$, hence $\gamma=0$.
If no such $\alpha_0$ exists, \eqref{eq:linear-LW-relation} directly gives
$\gamma=0$.  Thus every $d_\alpha$ is zero, so
$I_{q,\Delta}\equiv0$.

Now suppose $\deg q=n\ge2$.  In the notation of \eqref{eq:Uj}, one has
$U_0=I_{q,\Delta}$.  If $U_0(\xi)$ is algebraic, then
\cref{prop:value-to-exp-poly} gives the exponential-polynomial representation
\eqref{eq:exp-polynomial-U0}.  The annihilation argument of
\cref{sec:annihilation} produces a nonzero polynomial $H$ satisfying the
phase-moment identities \eqref{eq:phase-moments}.  Applying
\cref{thm:phase-moment-rigidity} yields $U_0\equiv0$.  This proves the
converse implication and the theorem.
\end{proof}

\begin{proof}[Proof of \cref{cor:single-integral}]
Take $\Delta=[B]-[A]$.  Since
\[
  I_{q,\Delta}(0)=\int_A^B1\,dx=B-A\ne0,
\]
the entire function $I_{q,\Delta}$ is not identically zero.  Apply
\cref{thm:algebraic-value-rigidity}.
\end{proof}

\begin{proof}[Proof of \cref{thm:relation-transfer}]
Let
\[
  \Delta=\sum_{j=1}^r a_j\Delta_j.
\]
Then $I_{q,\Delta}=\sum_{j=1}^r a_jI_j$.  If the left side of
\eqref{eq:relation-transfer} vanishes, then
\[
  I_{q,\Delta}(\xi)=-a_0\in\Qbar.
\]
By \cref{thm:algebraic-value-rigidity}, $I_{q,\Delta}\equiv0$ and its value at
$\xi$ is zero.  Therefore $a_0=0$.  The converse is immediate.  The dimension
identity \eqref{eq:dimension-transfer} is the corresponding equality of
relation spaces.
\end{proof}

\begin{lemma}[Separated regular fibers give functional independence]
\label{lem:regular-fiber-functional-independence}
Let $a_0,\ldots,a_r\in\C$ be such that
$q(a_0),\ldots,q(a_r)$ are pairwise distinct regular values of $q$.  If
\[
  \Delta=\sum_{i=0}^r d_i[a_i],
  \qquad \sum_{i=0}^r d_i=0,
\]
and $I_{q,\Delta}\equiv0$, then $d_0=\cdots=d_r=0$.
\end{lemma}

\begin{proof}
Choose the star in \cref{lem:constellation} so that every
$t_i=q(a_i)$ is the terminal endpoint of its own arm.  Because $t_i$ is a
regular value, every point of $q^{-1}(t_i)$ is unramified, and every lift of
the terminal arm has exactly one incident edge at its terminal vertex.
Orient the terminal arm toward $t_i$ and use a graph chain
$\widetilde\Gamma$ with boundary $\Delta$.

Let $x_{i,1},\ldots,x_{i,n}$ be the inverse branches over the open terminal
arm, numbered so that $x_{i,1}(t)\to a_i$ as $t\to t_i$.  If $f_{i,k}$ is the
coefficient of the corresponding lifted edge in $\widetilde\Gamma$, then the
boundary coefficient at its terminal vertex is exactly $f_{i,k}$.  Since the
only point of the support of $\Delta$ in the fiber over $t_i$ is $a_i$, one
has
\[
  f_{i,1}=d_i,
  \qquad
  f_{i,k}=0\quad(k\ge2).
\]
By the edge-density criterion \eqref{eq:edge-criterion}, the density on this
arm vanishes identically.  Hence
\[
  0=\psi_i(t)=\frac{d_i}{q'(x_{i,1}(t))}.
\]
The denominator is nonzero, so $d_i=0$.  This holds for every $i$.
\end{proof}

\begin{proof}[Proof of \cref{cor:regular-value-independence}]
Put
\[
  I_j(z)=\int_{a_0}^{a_j}e^{zq(x)}\,dx,
  \qquad 1\le j\le r.
\]
If $\sum_{j=1}^r c_jI_j(z)\equiv0$, then the corresponding boundary zero-chain
is
\[
  \sum_{j=1}^r c_j([a_j]-[a_0])
  =\sum_{j=1}^r c_j[a_j]
   -\left(\sum_{j=1}^r c_j\right)[a_0].
\]
By \cref{lem:regular-fiber-functional-independence}, all its coefficients
vanish, and therefore $c_1=\cdots=c_r=0$.  Thus $I_1,\ldots,I_r$ are
linearly independent over $\Qbar$ as entire functions.  Apply
\cref{thm:relation-transfer}.
\end{proof}

\section*{Acknowledgments and AI-assisted research disclosure}

This manuscript was developed through interactive work between Christopher
D. Long and ChatGPT 5.6 Pro.  ChatGPT 5.6 Pro assisted in proof discovery,
organization, symbolic checking, reference verification, adversarial auditing,
and drafting.  Claude Opus 5 contributed the complete-reducibility/projector strategy,
the reduction to phase-polynomial moments, and extensive adversarial audits.
The AI systems are not authors.  The human author bears full responsibility
for all statements, proofs, citations, and any remaining errors.  This draft
has not yet undergone independent peer review or formal verification.

\small
\begin{thebibliography}{99}

\bibitem{AdamczewskiRivoal}
B.~Adamczewski and T.~Rivoal,
\emph{Exceptional values of $E$-functions at algebraic points},
Bull. Lond. Math. Soc. \textbf{50} (2018), no.~4, 697--708.
\href{https://doi.org/10.1112/blms.12168}{doi:10.1112/blms.12168}.

\bibitem{Baker}
A.~Baker,
\emph{Transcendental Number Theory},
Cambridge Mathematical Library,
Cambridge University Press, Cambridge, 1990.

\bibitem{Beukers}
F.~Beukers,
\emph{A refined version of the Siegel--Shidlovskii theorem},
Ann. of Math. (2) \textbf{163} (2006), no.~1, 369--379.
\href{https://doi.org/10.4007/annals.2006.163.369}{doi:10.4007/annals.2006.163.369}.

\bibitem{BostanRivoalSalvy}
A.~Bostan, T.~Rivoal, and B.~Salvy,
\emph{Minimization of differential equations and algebraic values of
$E$-functions},
Math. Comp. \textbf{93} (2024), no.~347, 1427--1472.
\href{https://doi.org/10.1090/mcom/3912}{doi:10.1090/mcom/3912}.

\bibitem{deCataldoMigliorini}
M.~A.~A. de~Cataldo and L.~Migliorini,
\emph{The decomposition theorem, perverse sheaves and the topology of
algebraic maps},
Bull. Amer. Math. Soc. (N.S.) \textbf{46} (2009), no.~4, 535--633.
\href{https://doi.org/10.1090/S0273-0979-09-01260-9}{doi:10.1090/S0273-0979-09-01260-9}.

\bibitem{Delaygue}
\'E.~Delaygue,
\emph{A Lindemann--Weierstrass theorem for $E$-functions},
J. Reine Angew. Math. \textbf{820} (2025), 75--85.
\href{https://doi.org/10.1515/crelle-2024-0090}{doi:10.1515/crelle-2024-0090}.

\bibitem{HTT}
R.~Hotta, K.~Takeuchi, and T.~Tanisaki,
\emph{$D$-Modules, Perverse Sheaves, and Representation Theory},
Progress in Mathematics, vol.~236,
Birkh\"auser Boston, Boston, MA, 2008.
\href{https://doi.org/10.1007/978-0-8176-4523-6}{doi:10.1007/978-0-8176-4523-6}.

\bibitem{NesterenkoShidlovskii}
Yu.~V. Nesterenko and A.~B. Shidlovskii,
\emph{Linear independence of values of $E$-functions},
Mat. Sb. \textbf{187} (1996), no.~8, 93--108;
English transl., Sb. Math. \textbf{187} (1996), no.~8, 1197--1211.
\href{https://doi.org/10.1070/SM1996v187n08ABEH000152}{doi:10.1070/SM1996v187n08ABEH000152}.

\bibitem{PakovichMuzychuk}
F.~Pakovich and M.~Muzychuk,
\emph{Solution of the polynomial moment problem},
Proc. Lond. Math. Soc. (3) \textbf{99} (2009), no.~3, 633--657.
\href{https://doi.org/10.1112/plms/pdp010}{doi:10.1112/plms/pdp010}.

\bibitem{PakovichRoytvarfYomdin}
F.~Pakovich, N.~Roytvarf, and Y.~Yomdin,
\emph{Cauchy-type integrals of algebraic functions},
Israel J. Math. \textbf{144} (2004), 221--291.
\href{https://doi.org/10.1007/BF02916714}{doi:10.1007/BF02916714}.

\end{thebibliography}

\end{document}
